\section{模型评价与推广}

本节用一张结论与证据表收束全文，随后列出优势与不足各若干条。
每条结论都指向支撑它的图表编号，并注明适用边界。

\begin{table}[H]
\centering\small
\caption{主要结论、支撑证据与适用边界}
\label{tab:conclusions}
\begin{tabular}{p{4.2cm}p{3.4cm}p{4.6cm}}
\toprule
结论 & 支撑证据 & 适用边界 \\
\midrule
化学特征对品评小组均分有可复现的预测能力，主模型等级 MAE 0.479，显著优于只知风格的基线 0.527 & 表~\ref{tab:q1-main}、图~\ref{fig:q1-dim} & 连续 $R^2$ 只有 0.184，信度中位数 0.53，预测能力有限 \\
\addlinespace
只用附件无法还原原始等级，两阶段对照 C1 等级 MAE 0.786 & 表~\ref{tab:q1-main}、表~\ref{tab:prep-reference} & C1 按自身规则选定等级映射，仍以同一原始等级为真值 \\
\addlinespace
多数维度可靠支持的关键化合物是 1 到 2 个，第 3 位以报告为主 & 表~\ref{tab:q2-rep}、图~\ref{fig:q2-stability} & Top-3 是关联结果，不是因果主张 \\
\addlinespace
选择规则由模拟还原实验确定，只用化学矩阵，不依赖感官数据 & 图~\ref{fig:q2-recovery} & 模拟假定少数化合物按秩线性可加，与真实结构存在差异 \\
\addlinespace
酯类加标在独立实验中一致提高酯类评分，酒花与苦味对照不变 & 表~\ref{tab:q2-spike} & 混合物含 6 到 7 种化合物，不能归因到单一化合物 \\
\addlinespace
风格条件生成模型在留出啤酒上通过大部分验收项，缺陷约束在几乎所有风格对上起作用 & 表~\ref{tab:q3-gen}、图~\ref{fig:q3-fingerprint} & 验收阈值在预运行后固定，属工程标准，不是独立验证 \\
\addlinespace
迁移只沿一条线段搜索，$\alpha^*$ 小不说明不存在满足约束的目标风格谱 & 图~\ref{fig:q3-transfer-fig} & 结论针对选定路径与问题一模型 \\
\addlinespace
问题一主输出向中间等级收缩，几乎不给出 5 到 6 级，第二输出保持等级分布 & 表~\ref{tab:q1-rules}、图~\ref{fig:q1-grade} & 这是以等级 MAE 为主指标的直接后果 \\
\bottomrule
\end{tabular}
\end{table}

\subsection{优势}

问题一预测品评等级，问题二找出关键化合物，问题三生成风格条件化合物谱并迁移，三个问题各有一条可独立验证的路径。
问题一以逐品评员原始评分的啤酒均分为目标和真值，不依赖附件四从标准化值还原 0 到 6 分的换算；只用附件换算的对照 C1 端到端等级 MAE 为 0.786，主模型为 0.479，差距可直接读出。
问题一的主指标、主输出、成功标准与对照胜出时的处理都在正式运行前写定，正式运行后按声明执行，没有按结果切换模型或规则。
问题二的稳定性用日期组重抽量化，模拟还原实验只用化学矩阵，独立于感官数据，回应了根据结果调整规则的质疑。
问题三的迁移把题面的不恶化要求写成化学层与感官层两层可检验的约束，并在潜空间与原始尺度两处复查。

\subsection{不足}

化学数据只解释品评小组均分的一小部分方差，连续 $R^2$ 为 0.184，信度（同一批评分随机分两半后相关的外推值）中位数为 0.53，后者给出可解释方差的上限，两者并列说明预测能力接近测量噪声的上限。
问题一的主输出向中间收缩，几乎不给出 5 到 6 级，而真值中这两个等级占 0.968\%。
随机森林对照的等级 MAE 为 0.472，低于主模型的 0.479，但 QWK 为 0.251，低于主模型的 0.290；主模型是运行前选定的模型，不是外层比较中误差最低的模型。
问题二的第 3 位基本不稳定，入选频率按名次约为 0.60、0.36、0.22，与主要对照方法 M3 的 Top-1 一致比例只有 0.45，这是不同估计对象的差异，不是某一个错了。
问题三的生成样本比真实啤酒更接近风格均值，风格判别器对真实啤酒的识别准确率为 0.550，对生成样本为 0.481，生成样本围绕收缩后的风格均值分布。
问题三的感官层约束依赖问题一在缺陷维度上的预测，这些维度多为零分主导、信度低，化学层约束才是主要保证，感官层是模型层面的检查。
